\begin{corollary}[2.7.2]
\label{IV.2.7.2-fr}
Siendo las hipótesis y notaciones las de \textup{(2.7.1)}, supongamos $f$
cuasicompacto; sean $\mathcal L$ un $\mathcal O_X$-módulo invertible,
$\mathcal L'=\mathcal L\otimes_{\mathcal O_Y}\mathcal O_{Y'}$ su imagen
inversa. Para que $\mathcal L$ sea amplio \textup{(resp.} muy amplio\textup{)} con respecto a
$f$, es necesario y suficiente que $\mathcal L'$ sea amplio \textup{(resp.} muy
amplio\textup{)} con respecto a $f'$.
\end{corollary}

La condición es necesaria sin hipótesis sobre $g:S'\to S$ (\textbf{II}, 4.4.10 y
4.6.13); para ver que es suficiente, es posible, como en (2.7.1), limitarse al
caso en que $S=Y$, $S'=Y'$. La hipótesis sobre $\mathcal L'$ implica que
$f'$ es cuasicompacto y separado (\textbf{II}, 4.6.1), por lo que lo mismo ocurre con
$f$ ((2.6.2, (v)) y (2.7.1, (i))). Pongamos
$\mathcal E=f_*(\mathcal L)$, $\mathcal E'=f'_*(\mathcal L')$; se sigue de
(2.3.1) que el homomorfismo canónico
$u:g^*(\mathcal E)\to\mathcal E'$ es biyectivo. Si $\mathcal L'$ es \emph{muy
amplio} con respecto a $f$, el homomorfismo canónico
$\sigma':f'^*(\mathcal E')\to\mathcal L'$ es sobreyectivo, y el morfismo
$r'=r_{\mathcal L',\sigma'}:X'\to\mathbf P(\mathcal E')$
es una \emph{inmersión} (\textbf{II}, 4.4.4, \emph{b)}) necesariamente cuasicompacta (1.1.2,
(v)). El hecho de que $u:g^*(\mathcal E)\to\mathcal E'$ sea biyectivo
implica que, si $h:\mathbf P(\mathcal E)\to Y$,
$h':\mathbf P(\mathcal E')\to Y'$ son los morfismos estructurales, $h'$ se
identifica con $h_{(Y')}$ (\textbf{II}, 4.1.3); por otra parte, denotando por $g'$ la
proyección $X'\to X$, $g'$ es fielmente plano (2.2.13), se tiene
$f\circ g'=g\circ f'$, y se comprueba fácilmente que el homomorfismo
$g'^*(\sigma):g'^*(f^*(f_*(\mathcal L)))\to g'^*(\mathcal L)$
es idéntico al homomorfismo compuesto
\[
f'^*(g^*(f_*(\mathcal L)))\xrightarrow{f'^*(u)}
f'^*(f'_*(\mathcal L'))\xrightarrow{\sigma'}\mathcal L'
\]
(por ejemplo, reduciéndose al caso en que $Y$ y $Y'$ son afines). Como
$\sigma'$ es sobreyectivo y $f'^*(u)$ es biyectivo, se ve que
$g'^*(\sigma)$ es sobreyectivo, por lo que lo mismo ocurre con $\sigma$
(2.2.7). Se concluye que el morfismo
$r=r_{\mathcal L,\sigma}:X\to\mathbf P(\mathcal E)$ está definido en todas
partes (\textbf{II}, 3.7.4); además, si se pone $P=\mathbf P(\mathcal E)$,
$P'=\mathbf P(\mathcal E')$ y si $g''$ es la proyección $P'\to P$, $r'$ se
identifica con $r_{(P')}$ (\textbf{II}, 4.2.10), y $g''$ es fielmente plano y
cuasicompacto (1.1.2 y 2.2.13). Se concluye, pues, de (2.7.1, (xi)) que $r$
es una inmersión y, por consiguiente, $\mathcal L$ es \emph{muy amplio} (\textbf{II}, 4.4.4,
\emph{b)}).

Supongamos ahora que $\mathcal L'$ es \emph{amplio} con respecto a $f'$; para
demostrar que $\mathcal L$ es amplio con respecto a $f$, es posible limitarse
al caso en que $Y$ es afín (\textbf{II}, 4.6.4), y, en virtud de (2.2.12) y de (\textbf{II},
4.6.13), se puede suponer también que $Y'$ es afín. Entonces $X$ y $X'$ son

\oldpage[IV-2]{33}
esquemas cuasicompactos, y, para demostrar que $\mathcal L$ es $f$-amplio,
se puede aplicar el criterio de (\textbf{II}, 4.6.8, \emph{c)}). Sea, pues, $\mathcal F$ un
$\mathcal O_X$-módulo cuasicoherente de tipo finito; si
$\sigma:f^*(f_*(\mathcal F\otimes\mathcal L^{\otimes n}))
\to\mathcal F\otimes\mathcal L^{\otimes n}$
es el homomorfismo canónico, se ve, como antes, que $g'^*(\sigma)$ es el
homomorfismo compuesto
\[
f'^*(g^*(f_*(\mathcal F\otimes\mathcal L^{\otimes n})))
\xrightarrow{f'^*(u)}f'^*(f'_*(\mathcal F'\otimes
\mathcal L'^{\otimes n}))\xrightarrow{\sigma'}
\mathcal F'\otimes\mathcal L'^{\otimes n}
\]
poniendo $\mathcal F'=g'^*(\mathcal F)$, teniendo en cuenta
$(\mathbf{0}_{\mathrm I},4.3.3.1)$ y denotando por $u$ el homomorfismo canónico
$g^*(f_*(\mathcal F\otimes\mathcal L^{\otimes n}))
\to f'_*(\mathcal F'\otimes\mathcal L'^{\otimes n})$.
Ahora bien, se sabe que $u$ es biyectivo cualquiera que sea $n$ (2.3.1); por
otra parte, como $\mathcal F'$ es cuasicoherente y de tipo finito, la
hipótesis de que $\mathcal L'$ es amplio con respecto a $f'$ implica la
existencia de un $n_0$ tal que para $n\geq n_0$, $\sigma'$ sea sobreyectivo;
se ve, pues, que $g'^*(\sigma)$ es sobreyectivo para $n\geq n_0$, y, como
$g'$ es fielmente plano, $\sigma$ es sobreyectivo para estos valores de $n$
(2.2.7), lo que termina la demostración.
